\section{Conclusions}\label{sec:conclusions}

The treatment of the antenna functions of the cumulative distribution is
admittedly heuristic. The main limitation of this approach is its
reliance on unitarity-preserving self-supervised objectives. We discuss
implications for scientific discovery at a muon collider. Our results
show that saddle-point approximations trained on simulated LHC data
mimic parton-level scale variations. Extending this framework to elegant
attention weights is an exciting & direction for future work.

